% New English translation from checked cote 25, site 106 (PDF page 107).
% AI-translated by GPT-6 Astra (OpenAI, ChatGPT Pro mode).
\scanpage{25/106}
\mspage{9}

2. \uline{Relative divisors and invertible sheaves on projective}

\uline{and multiprojective bundles}.

\hand{2.1.} As in no.1, let P=P($\underline{E}$), $\underline{E}$ loc free on S of rank $\hand{\geqslant}$ 2 everywhere.
We propose to determine the set Div(P/S) of relative divisors
\var{\inserted{$\geqslant 0$}}{$\geqslant 0$}{insertion}{ega25_25_106-v0001} on P w.r.t. S. We know that the datum of such a divisor amounts to the
datum of an invertible Module $\underline{L}$ on P, equipped with a transversely regular
section $\hand{\phi}$ of $\underline{L}$. But by 1.1. (disregarding a
possible partition of S, when S is not connected) $\underline{L}$ is isomorphic
to an $\underline{M}\otimes\underline{O}_P(n)$, where $\underline{M}$ is an invertible Module on S, moreover determined
up to unique isomorphism in terms of $\underline{L}$. Moreover, we know
that we then have canonical isoms
\[ \hand{(*)} \qquad f_{\hand{*}}(\underline{L}) \hand{\simeq} \underline{M}\otimes f_{\hand{*}}(\underline{O}_P(\unsure{n|m})) \hand{\simeq} \underline{M}\otimes\underline{Sym}^n(\underline{E}) \]
so the datum of a section \var{\xout{of}}{}{typewriter strike-out}{ega25_25_106-v0002} $\hand{\phi}$ of $\underline{E}$ amounts to the datum of a section
\var{\inserted{$\phi$}}{$\phi$}{insertion}{ega25_25_106-v0003} of $\underline{M}\otimes\underline{Sym}^n$(E). Taking account of the fact that the fibers of P/S are integral,
we moreover see that $\hand{\phi}$ is transversely regular (i.e. regular
on every fiber) iff \hand{$\phi$}(s)$\hand{\neq}$0 for every s$\hand{\in}$S, or equivalently,
i\hand{ff} the homomorphism
\[ \hand{\phi} : \overset{\hand{\vee}}{\underline{M}} \hand{\longrightarrow} \underline{Sym}^n(E) \]
defined by \hand{$\phi$} is ``universally injective'' i.e. loc an isomorphism onto
a direct factor, or equivalently, if its transpose
\[ \hand{\check{\phi}} : \underline{Sym}^n(\underline{E})^{\hand{\vee}} \hand{\longrightarrow} \underline{M} \]
is surjective.

We will say that the relative divisor D on P is \uline{of degree} n
if $\underline{O}_X(\underline{D})= \underline{L}$ is of degree n in the sense of the preceding no. \var{\deleted{Thus, by 1.1.}}{}{deletion}{ega25_25_106-v0004}
\var{\deleted{for given D, there exists a unique \var{\inserted[typed]{decomposition}}{decomposition}{insertion}{ega25_25_106-v0006} \var{\xout{partition}}{}{typewriter strike-out}{ega25_25_106-v0007} of S as a sum of opens $S_n$}}{}{deletion}{ega25_25_106-v0005}
\var{\deleted{\unsure{(n N)|n N)} such}}{}{deletion}{ega25_25_106-v0008}
\tnote{25-106-1}{The abbreviations p. r. à and sss are translated as w.r.t. and iff. In the second iff, the handwritten ff corresponds to the source's handwritten ss. The source says a section of E after the first formula and uses O with subscript X in the degree definition; both are retained.}

\mspage{106}
